\label{sec:validation}

This section separates validation and external-comparison evidence from the
performance results that follow. The distinction is important because the
paper uses different levels of evidence for different claims: Hyperelasticity
has a narrow external companion comparison against JAX-FEM, while
\plasticitythree{} has a source-family validation ladder for the implemented
endpoint surrogate.

\subsection{Hyperelasticity Companion Comparison}
\label{sec:val-hyperelasticity}

The hyperelasticity problem definition, repository energy, and load path are
defined in \Cref{sec:bench-he,eq:hyperelasticity-discrete}. The external
comparison here is deliberately narrow. It uses a serial companion case against
JAX-FEM on the same mesh and prescribed displacement schedule, then
post-compares both terminal states with the repository energy from
\cref{eq:hyperelasticity-energy}. This supports a solution-level
implementation comparison for this hyperelastic problem only. It should not be
read as a plasticity validation, a broad framework ranking, or constitutive-law
identity between the two implementations: the JAX-FEM worker uses a
logarithmic-$J$ Piola form, whereas the repository energy uses the
$D_1(J-1)^2$ volumetric penalty in \cref{eq:hyperelasticity-energy}.

\begin{figure}[H]
  \centering
  \begin{subfigure}[t]{0.46\PaperFigureFull}
    \centering
    \includegraphics[width=\linewidth]{jax_fem_hyperelastic_baseline_energy_history.pdf}
    \caption{Terminal energy over the matched displacement schedule.}
  \end{subfigure}
  \hfill
  \begin{subfigure}[t]{0.46\PaperFigureFull}
    \centering
    \includegraphics[width=\linewidth]{jax_fem_hyperelastic_baseline_centerline.pdf}
    \caption{Final centerline $u_x$ profile on the common sample path.}
  \end{subfigure}
  \caption{Serial hyperelastic companion comparison against JAX-FEM on the
  same mesh and prescribed displacement schedule, with repository-energy
  post-comparison of the terminal states.}
  \label{fig:hyperelasticity-jaxfem-baseline}
\end{figure}

\begin{table}[H]
  \centering
  \begin{tabular*}{\textwidth}{@{}>{\RaggedRight\arraybackslash}p{0.15\textwidth}@{\hspace{1.0em}}>{\RaggedRight\arraybackslash}p{0.30\textwidth}@{\extracolsep{\fill}}c c c@{}}
\toprule
Group & Quantity & Relative difference & Median wall time [s] & Status \\
\midrule
Agreement & final energy & $\num{8.24}\times 10^{-6}$ & -- & -- \\
Agreement & full-field displacement relative $L^2$ & $\num{2.01}\times 10^{-4}$ & -- & -- \\
Agreement & centerline relative $L^2$ & $\num{5.47}\times 10^{-5}$ & -- & -- \\
Agreement & $u_{\max}$ curve relative $L^2$ & $\num{8.85}\times 10^{-16}$ & -- & -- \\
\addlinespace
Timing & this work serial direct & -- & \num{0.176} & -- \\
Timing & JAX-FEM UMFPACK serial & -- & \num{2.52} & -- \\
\addlinespace
Contract/gate & same mesh path & -- & -- & pass \\
Contract/gate & same displacement schedule & -- & -- & pass \\
Contract/gate & agreement thresholds & -- & -- & pass \\
Contract/gate & repository-energy post-comparison & -- & -- & used \\
\bottomrule
\end{tabular*}

  \caption{Fairness-gated hyperelastic companion comparison against JAX-FEM.}
  \label{tab:hyperelasticity-jaxfem-baseline}
\end{table}

The fairness gate in \Cref{tab:hyperelasticity-jaxfem-baseline} passes on all
checked quantities. The final energy relative difference is
\num{8.244e-6}, the full-field displacement relative $L^2$ difference is
\num{2.011e-4}, the centerline relative $L^2$ difference is
\num{5.473e-5}, and the $u_{\max}$-curve relative $L^2$ difference is
effectively zero. On this small serial CPU companion case, the maintained
direct solve is faster (\SI{0.176}{s} median wall time versus \SI{2.524}{s}
for JAX-FEM), but the scientific role of this table is solution agreement
under a narrow hyperelastic comparison contract rather than framework ranking.

\subsection{Plasticity3D Validation Ladder}
\label{sec:val-plasticity3d}

The \plasticitythree{} benchmark is tied to a Mohr--Coulomb slope-stability
source family, but the repository object studied here is the algorithmic
endpoint surrogate defined in
\Cref{sec:bench-plasticity3d,eq:plasticity3d-discrete,eq:plasticity3d-potential}.
The Sysala return-mapping, optimization, convex-optimization, and advanced
continuation papers provide source-family and method context for this problem
class \citep{sysala2017returnmapping,sysala2021optimization,sysala2025convexoptimization,sysala2025advancedcontinuation};
they are not used here as verified numeric paper baselines.

The validation ladder therefore separates two questions. Layer 1A tests
same-case source-faithfulness of the endpoint surrogate on the direct-branch
$P_2(L_1)$ glued-bottom case. This layer is tight: the work relative difference
is \num{3.877e-5}, the displacement relative $L^2$ difference is
\num{3.517e-3}, and the deviatoric-strain relative $L^2$ difference is
\num{8.720e-3}. Layer 2 is different: it is a fixed-$\lambda$ diagnostic
against a source-operator comparison, not proof of path-consistent
incremental-history equivalence. After matching the source boundary contract
to the maintained glued-bottom free-DOF mask, the validation schedule is capped
at the main-figure load level $\lambda_{\mathrm{sr}}=1.55$. It keeps the
highest-successful-$\lambda$ schedule proxy aligned on the tested grid and
passes the configured gates for the $u_{\max}(\lambda)$ curve and endpoint
displacement: the $u_{\max}(\lambda)$ relative $L^2$ difference is
\num{6.396e-6}, and the endpoint displacement relative $L^2$ difference is
\num{8.299e-4}. The endpoint deviatoric-strain relative $L^2$ difference
(\num{4.307e-3}) and boundary-profile relative $L^2$ difference
(\num{5.229e-4}) are reported as diagnostics rather than acceptance gates.
This comparison uses a matched mesh, material table, Davis-B strength
reduction, load schedule, and glued-bottom free-DOF mask, but it remains a
fixed-load endpoint check rather than a path-consistent plastic-history
reference.

\begin{figure}[H]
  \centering
  \begin{subfigure}[t]{0.46\PaperFigureFull}
    \centering
    \includegraphics[width=\linewidth]{plasticity3d_validation_layer1a_boundary.pdf}
    \caption{Layer 1A: same-case source-faithfulness on the surrogate.}
  \end{subfigure}
  \hfill
  \begin{subfigure}[t]{0.46\PaperFigureFull}
    \centering
    \includegraphics[width=\linewidth]{plasticity3d_validation_umax_curve.pdf}
    \caption{Layer 2: fixed-$\lambda$ source-operator diagnostic.}
  \end{subfigure}
  \caption{Plasticity3D validation ladder. The left panel documents tight
  same-case source-faithfulness for the implemented endpoint surrogate; the
  right panel shows close fixed-$\lambda$ source-operator agreement after
  matching the glued-bottom boundary contract.}
  \label{fig:plasticity3d-validation}
\end{figure}

\begin{table}[H]
  \centering
  \begin{tabular*}{\textwidth}{@{}c@{\hspace{1.0em}}>{\RaggedRight\arraybackslash}p{0.36\textwidth}@{\extracolsep{\fill}}c c@{}}
\toprule
Layer & Comparison & Relative difference & Status \\
\midrule
1A & work & $\num{3.88}\times 10^{-5}$ & -- \\
1A & displacement relative $L^2$ & $\num{3.52}\times 10^{-3}$ & -- \\
1A & deviatoric-strain relative $L^2$ & $\num{8.72}\times 10^{-3}$ & -- \\
2 & highest-successful $\lambda$ & \num{0} & pass \\
2 & $u_{\max}(\lambda)$ relative $L^2$ & $\num{6.4}\times 10^{-6}$ & pass \\
2 & endpoint displacement relative $L^2$ & $\num{8.3}\times 10^{-4}$ & pass \\
2 & endpoint deviatoric-strain relative $L^2$ & $\num{4.31}\times 10^{-3}$ & diagnostic \\
2 & boundary profile relative $L^2$ & $\num{5.23}\times 10^{-4}$ & diagnostic \\
2 & acceptance gate & -- & pass \\
\bottomrule
\end{tabular*}

  \caption{Plasticity3D validation summary for same-case source-faithfulness
  (layer 1A) and the fixed-$\lambda$ source-operator diagnostic (layer 2).}
  \label{tab:plasticity3d-validation}
\end{table}

The resulting interpretation is limited but useful. \plasticitythree{} is
validated here as a source-faithful maintained endpoint-surrogate workflow in
Layer 1A, while Layer 2 supports the same endpoint surrogate under a fixed-load
source-operator contract. Neither layer is a path-consistent
incremental-plasticity history validation.
